\section{Effective convergence and differentiation}\label{sec:convergence}

\subsection{A primitive recursive ratio cutoff}

Fix complex-rational parameters \(a,b,c\) and a rational radius
\(0\le r<1\).  Let rational numbers \(A,B,C\) satisfy
\[
 \norm a\le A,\qquad \norm b\le B,\qquad \norm c\le C.
\]
For \(n>C\), the reverse triangle inequality gives
\(\norm{n+c}\ge n-C\).  Hence, for \(\norm z\le r\),
\begin{equation}\label{eq:term-ratio-bound}
 \norm{\frac{A_{n+1}z^{n+1}}{A_nz^n}}
 \le r\frac{(n+A)(n+B)}{(n+1)(n-C)}
 \le r\frac{(n+A)(n+B)}{n(n-C)}.
\end{equation}
This estimate is used only after denominator apartness has been verified.

\begin{lemma}[Explicit ratio cutoff]\label{lem:ratio-cutoff}
If \(0<r<1\), put
\[
 q=\frac{1+r}{2},\qquad
 \varepsilon=\frac{q-r}{r}=\frac{1-r}{2r},
\]
\[
 D=(1+\varepsilon)C+A+B+AB,
 \qquad
 N_0=1+\left\lceil\max\left(C,\frac{2D}{\varepsilon}\right)\right\rceil.
\]
Then for every \(n\ge N_0\), the right side of
\eqref{eq:term-ratio-bound} is at most \(q\).  For \(r=0\), every nonconstant
term vanishes.
\end{lemma}

\begin{proof}
It is enough to prove
\[
 (n+A)(n+B)\le(1+\varepsilon)n(n-C).
\]
After expansion this is
\[
 \varepsilon n^2\ge
  ((1+\varepsilon)C+A+B)n+AB.
\]
For \(n\ge1\), the right side is at most \(Dn\).  The definition of
\(N_0\) gives \(\varepsilon n\ge2D\), hence the desired inequality with
room to spare.  Every comparison is rational and every ceiling is obtained by
bounded integer search.
\end{proof}

\begin{theorem}[Uniform convergence on a strict subdisk]\label{thm:disk-convergence}
Assume the denominator factors \(c+j\) carry apartness certificates.  For
every rational \(r<1\), the partial sums
\(H_N(a,b;c;z)\) have a primitive recursive Cauchy modulus uniformly on
\(\norm z\le r\).  More precisely, for \(N\ge N_0\) and \(m\ge0\),
\begin{equation}\label{eq:geometric-tail}
 \norm{\sum_{n\ge N+m}A_nz^n}
 \le \frac{\norm{A_N}r^N}{1-q}q^m.
\end{equation}
A primitive recursive bounded search finds the least \(m\) making the right
side at most \(2^{-p}\).
\end{theorem}

\begin{proof}
By \cref{lem:ratio-cutoff}, every term after \(N_0\) is at most \(q\) times
the preceding term.  Induction gives
\(\norm{A_{N+j}z^{N+j}}\le\norm{A_N}r^Nq^j\).
The finite geometric identity
\[
 (1-q)\sum_{j<m}q^j=1-q^m
\]
is an exact polynomial identity, so the tail bound follows by passing to the
Cauchy estimate.  No classical ratio test is invoked: the ratio cutoff and
the geometric remainder are explicit.

For termination of the search, write \(d=1-q>0\).  The finite binomial
inequality
\((1-d)^{-m}\ge1+md\) gives
\(q^m\le(1+md)^{-1}\).  Taking
\(m\ge2^{p+s}/d\), where \(s\) absorbs
\(\norm{A_N}r^N/(1-q)\), certainly succeeds.  This is a primitive recursive,
though intentionally crude, bound.
\end{proof}

\begin{corollary}[Named analytic function on a strict subdisk]
For each fixed guarded parameter package and rational \(r<1\),
\(z\mapsto\F a b c z\) is a named uniformly continuous function on
\(\norm z\le r\).  An evaluator computes a sufficiently long rational
partial sum; a modulus follows by splitting
\[
 F(z)-F(w)=\sum_{n<N}A_n(z^n-w^n)+\Tail_N(z)-\Tail_N(w).
\]
\end{corollary}

\begin{proof}
Choose \(N\) so that both tails are at most \(2^{-p-2}\).  The finite head is
uniformly continuous because
\[
 z^n-w^n=(z-w)\sum_{j<n}z^jw^{n-1-j}.
\]
On \(\norm z,\norm w\le r\), its norm is at most
\(n r^{n-1}\norm{z-w}\).  The finite sum of these rational constants gives an
explicit Lipschitz bound and hence a dyadic modulus.
\end{proof}

\subsection{Analytic differentiation}

The formal identity of \cref{thm:derivatives} becomes an analytic identity
once both sides have compatible moduli.

\begin{theorem}[Effective termwise differentiation]\label{thm:analytic-derivative}
Let \(0\le r<\rho<1\) be rational.  On \(\norm z\le r\), the difference
quotient of the named sum converges effectively to the named derivative
series, and
\[
 \frac{d^s}{dz^s}\F a b c z
 =\frac{\poch a s\poch b s}{\poch c s}
   \F{a+s}{b+s}{c+s}{z}.
\]
after the shifted denominator guards are supplied.
\end{theorem}

\begin{proof}
Choose the series on the larger disk \(\norm z\le\rho\).  Its tail has a
primitive recursive modulus by \cref{thm:disk-convergence}.  For a monomial,
\[
 \frac{(z+h)^n-z^n}{h}-nz^{n-1}
 =\sum_{j=2}^{n}\binom nj z^{n-j}h^{j-1}.
\]
On \(\norm z\le r\) and \(\norm h\le(\rho-r)/2\), the norm is bounded by an
explicit finite multiple of \(\norm h\rho^{n-2}\).  Split the derivative
series into a finite head and tail.  Choose the tail cutoff first, then choose
\(h\) small enough for the finite head.  All choices are bounded searches over
rational inequalities.  The coefficient identity from
\cref{thm:derivatives} identifies the resulting derivative name.  Iterating
this argument gives fixed order \(s\); joint primitive recursive data give a
uniform variable-order schema.
\end{proof}

\subsection{Convergence on the unit circle}

The source pages state convergence at \(|z|=1\) when
\(\RePart(c-a-b)>0\).  A constructive version needs a positive margin, not
merely the negation of a nonpositive real part.

\begin{contract}[Boundary convergence data]\label{con:boundary}
Supply a positive rational \(\delta\) with
\[
 \RePart(c-a-b)\ge\delta>0,
\]
a denominator-apartness table, and an effective gamma-ratio estimate
\begin{equation}\label{eq:gamma-ratio-contract}
 \norm{A_n(a,b;c)}\le K(n+1)^{-1-\delta}
 \qquad(n\ge N_1)
\end{equation}
for rational \(K,N_1\).  The gamma/Stirling package in the repository gives a
route to such an estimate for fixed primitive recursive parameter names on a
compact nonresonant set.
\end{contract}

\begin{theorem}[Effective unit-circle convergence]\label{thm:unit-circle}
Under \cref{con:boundary}, the hypergeometric series converges absolutely and
uniformly on the full coded Euclidean unit disk
\[
   (\RePart z)^2+(\ImPart z)^2\le1,
\]
with a primitive recursive tail modulus.  In particular it converges
uniformly on the rational diamond \(\norm z\le1\).
\end{theorem}

\begin{proof}
Choose an integer \(k\ge1\) with \(1/k\le\delta\).  For \(N\ge N_1\), group
indices into blocks
\([j^k,(j+1)^k)\).  Every term in the block is at most
\(Kj^{-k-1}\), while the block has at most
\(k(j+1)^{k-1}\) terms.  Hence its contribution is at most
\(K'j^{-2}\) for an explicit rational \(K'\).  The elementary telescoping
bound
\[
 \sum_{j\ge J}\frac1{j^2}
 \le\sum_{j\ge J}\frac1{j(j-1)}=\frac1{J-1}
\]
gives a primitive recursive tail cutoff.  This avoids invoking an integral
test or a completed real-number order.  If
\((\RePart z)^2+(\ImPart z)^2\le1\), then the same inequality, raised to the
\(n\)-th power, certifies that each real and imaginary component of \(z^n\)
has absolute value at most one.  Consequently
\(\norm{A_nz^n}\le2\norm{A_n}\); the harmless factor two is absorbed in
\(K'\).  On \(\norm z\le1\), even that extra argument is unnecessary.
\end{proof}

\begin{remark}[Endpoint without a margin]
The qualitative assertion \(\RePart(c-a-b)>0\) is read as the existence of a
positive rational witness \(\delta\).  For a fixed named real with no supplied
witness, iPRA cannot in general decide on which side of zero it lies.  The
margin is mathematical data, not bureaucratic overhead.
\end{remark}

\subsection{Coefficientwise confluent limit}

Let \(h\ne0\) and consider
\[
 G_h(z)=\F a{h^{-1}}c{hz}.
\]
A coefficient contains
\[
 \poch{h^{-1}}n h^n=\prod_{j<n}(1+jh).
\]

\begin{theorem}[Effective confluent limit]\label{thm:confluent}
For every fixed \(n\),
\[
 \prod_{j<n}(1+jh)\longrightarrow1
\]
with the explicit bound
\begin{equation}\label{eq:finite-product-perturb}
 \norm{\prod_{j<n}(1+jh)-1}
 \le \left(1+n\norm h\right)^n-1.
\end{equation}
Consequently the coefficients of \(G_h\) converge to those of
\[
 {}_1F_1(a;c;z)=\sum_{n\ge0}\frac{\poch a n}{\poch c n}\frac{z^n}{n!}.
\]
On each rational disk \(\norm z\le r\), and for \(h\) restricted by a
sufficiently small rational bound, the convergence is uniform with a primitive
recursive modulus.
\end{theorem}

\begin{proof}
Expand the difference of two finite products by telescoping one factor at a
time:
\[
 \prod_{j<n}(1+u_j)-1
 =\sum_{k<n}u_k\prod_{j<k}(1+u_j).
\]
With \(u_j=jh\), a coarse estimate gives
\eqref{eq:finite-product-perturb}.  For uniform convergence, choose a head
cutoff \(N\) using a common ratio majorant for all sufficiently small
\(\norm h\).  On the finite head, choose \(h\) so that each coefficient error
is below the allocated budget.  The two tails use the same geometric bound.
This is the constructive head--tail proof; no interchange of two abstract
limits is needed.
\end{proof}
